Inverse kinematics (IK) is the problem of finding control parameters that bring a robot into a desired pose within its kinematic constraints. For manipulators of six or fewer joints, this problem is practically solvable and has at most 16 different solutions \cite{Manocha}. Adding additional joints results in an under-constrained problem with infinitely many feasible solutions satisfying the pose constraints. Therefore, it is interesting not simply to look for feasible solutions but optimal ones. In this paper, we focus on minimizing the total movement of joints. Specifically, the task is as follows: Given the Denavit-Hartenberg parameters \cite{Spong2020-mt} of a kinematic chain, input joint angles, and a desired pose for its end effector, find joint angles that achieve the desired pose while minimizing the weighted sum of differences from the input (preferred) angles.

The presented problem is very hard in general. The current state-of-the-art technique is based on \emph{sum of squares} (SOS) optimization and Gröbner bases, which are both NP-hard in the context of multivariate polynomial equations. Trutman \textit{et al.} demonstrate \cite{trutman2022globally} that for manipulators with at most seven joints, the optimal solutions admit a quadratic SOS representation, thereby being practically solvable. Such results do not extend to manipulators with more than seven joints.

A common way to avoid the complexity of IK is either to find an approximate solution or to design the robotic arm to be simple to control. 
Approaches from the first category include an approximation
based on the Jacobian by Buss \cite{Buss_undated-fo}, which uses numerical methods to find local optima. Similarly, Dai \textit{et al.} \cite{Dai2019-uf} designed a global solution method for the convex approximation of the original IK. On the other hand, the design of the \href{https://www.kuka.com/en-us/products/robotics-systems/industrial-robots/lbr-iiwa}{KUKA LBR iiwa} manipulator falls into the second category. However, having to design around computational constraints has the downside of limiting the achievable performance of a manipulator.

The problem may be further relaxed for real-time control by dropping the guarantees on optimality or convergence to the desired pose. For example, \href{https://robotology.github.io/robotology-documentation/doc/html/group__iKinSlv.html}{iKinSlv} finds a locally optimal solution which minimizes the distance from the desired pose using interior point optimization; the closed-loop inverse kinematics algorithm suggested in \href{https://github.com/stack-of-tasks/pinocchio/blob/97a00a983e66fc0a58667f9671e2d806ba9b730b/examples/inverse-kinematics.cpp}{the Pinocchio library} attempts to iteratively reduce the pose error, while modern approaches \cite{ExampleQP} based on convex \emph{quadratic programming} (QP) rely on asymptotic convergence.

Global solvers recover a globally optimal solution whenever it exists or guarantee that none exists. In the context of mechanical design, it is possible to optimize a design by computing and comparing the actual limits of performance in some desired metric specified by the user. Using the optimality guarantees to test existing solvers' properties is also possible. If a solver fails to provide a solution, it is possible to answer whether the instance was solvable.

The contribution of this paper is based on the premise that while SOS optimization is, in some sense, the correct answer to the \emph{polynomial optimization problem} (POP), incumbent implementations fall significantly short of their promised potential. We will show that highly optimized off-the-shelf global optimizers can converge faster for problems such as IK despite having much higher computational complexity. Instead of reformulating IK into a theoretically easier problem, we reformulate it into a problem for which efficient solvers exist.

Our approach finds globally optimal solutions of generic seven \emph{degree-of-freedom} (DOF) IK instances faster than the state-of-the-art. Our technique is unique in finding globally optimal IK solutions for generic manipulators with eight or more degrees of freedom.
